\section{Study B: Passenger Evaluation}
\label{sec:study-b}

The planned study uses a within-participant, repeated-measures design with three conditions: No Cue, Static Visual Control, and Synchronized Cue. Condition order will be balanced with a Latin square. The target is 30 completed participants, preceded by a 6--10 participant pilot. Participants will be adults in their early to mid twenties, with inclusion, exclusion, stopping, and missing-data rules finalized before recruitment.

Each condition is planned for a separate visit to reduce residual sickness and adaptation. The same vehicle, driver, route, time band, right-rear passenger seat, Windows laptop, and iPhone~16~Pro sensor companion will be used. Participants will complete matched reading-comprehension tasks. A small prompt will request an FMS score from 0 to 20 every 60~seconds under all conditions.

The primary outcome is the temporal development of FMS, summarized with area under the curve, peak value, and slope. Secondary outcomes include MSAQ change, task accuracy and response time, workload, cue distraction, preference, and withdrawal. The planned primary model is:
\begin{equation}
\begin{aligned}
\mathrm{FMS} \sim {} & \mathrm{Condition} \times \mathrm{Time} + \mathrm{MSSQ} \\
                      & + \mathrm{Order} + (1 \mid \mathrm{Participant}).
\end{aligned}
\end{equation}

Research-review status is currently \textbf{TBD}. No participant recruitment or data collection will begin until consent, safety procedures, and the applicable institutional review or exemption path have been resolved.
